% ============================================================
% common/preamble.tex —— 拆分版试卷项目 · 公共导言区
% 引擎：XeLaTeX；文档类 ctexart 由各版本壳文件自行声明
% 由合并书 math-gaokao-2017-2026/preamble.tex 改造而来：
%   - 页眉左右改为可重定义宏 \gkheadL / \gkheadR（各版本壳设置）
%   - hyperref 不再写死 pdftitle（壳内 \hypersetup 覆盖）
%   - 删除书籍专用的 \chaptermark 重定义与 \gkpaper 装配宏
% ============================================================

% ---------- 数学 ----------
\usepackage{amsmath}
\usepackage{amssymb}
\usepackage{mathtools}

% ---------- 绘图 ----------
\usepackage{tikz}
\usetikzlibrary{shapes.geometric, arrows.meta, calc, patterns, angles, quotes, positioning}
\usepackage{pgfplots}
\pgfplotsset{compat=1.18}

% ---------- 版面 ----------
\usepackage{geometry}
\usepackage{needspace}
\geometry{a4paper, margin=2.2cm, headheight=15pt}
% 紧凑版式（如印刷版）可在壳文件中于 \begin{document} 前再次调用 \geometry 覆盖
\usepackage{enumitem}
\setlist{itemsep=2pt, topsep=4pt, parsep=0pt}

% 中文与行内公式混排的断行弹性：
% emergencystretch 允许段落末路不通时追加伸展量；
% 降低二元/关系符罚分，允许行内长公式在 = 、+ 等处断行（否则整段外溢）
\emergencystretch=2em
\binoppenalty=700
\relpenalty=500
\usepackage{booktabs}      % 三线表
\usepackage[most]{tcolorbox}

% ---------- 选择题选项对齐 ----------
% tasks 宏包：\begin{gkopts}(列数) \task ... \end{gkopts}
% 自动等宽分栏，选项天然对齐；列数已由内容片段按选项长度写定（4/2/1）
\usepackage{tasks}
\settasks{label={}, label-width=0pt, item-indent=1.2em,
          column-sep=1em, after-item-skip=2pt, before-skip=2pt, after-skip=4pt}
\NewTasksEnvironment[style=itemize]{gkopts}[\task](4)

% ---------- 页眉页脚 ----------
\newcommand{\gkheadL}{}   % 左页眉：卷名（壳文件 \renewcommand 设置）
\newcommand{\gkheadR}{}   % 右页眉：版本名（壳文件 \renewcommand 设置）
\usepackage{fancyhdr}
\pagestyle{fancy}
\fancyhf{}
\fancyhead[L]{\small\gkheadL}
\fancyhead[R]{\small\gkheadR}
\fancyfoot[C]{\small —— 第 \thepage\ 页 ——}
\renewcommand{\headrulewidth}{0.4pt}

% ---------- 超链接（须最后加载）----------
\usepackage{hyperref}
\hypersetup{colorlinks=true, linkcolor=blue!55!black,
            bookmarksnumbered=true, unicode=true}
